\documentclass[10pt, letterpaper]{article}

\usepackage[T1]{fontenc}
\usepackage{lmodern}
\usepackage[margin=0.75in, top=0.70in, bottom=0.70in]{geometry}
\usepackage{amsmath, amssymb}
\usepackage{booktabs}
\usepackage{hyperref}
\usepackage{xcolor}
\usepackage{graphicx}
\usepackage{float}

\setlength{\parskip}{4pt plus 1pt minus 1pt}
\setlength{\parindent}{0pt}
\setlength{\textfloatsep}{8pt}
\setlength{\floatsep}{6pt}
\setlength{\intextsep}{6pt}
\setlength{\abovecaptionskip}{3pt}
\setlength{\belowcaptionskip}{2pt}

\let\olditemize\itemize
\renewcommand{\itemize}{\olditemize\setlength{\itemsep}{1pt}\setlength{\topsep}{2pt}}
\let\oldenumerate\enumerate
\renewcommand{\enumerate}{\oldenumerate\setlength{\itemsep}{1pt}\setlength{\topsep}{2pt}}

\hypersetup{colorlinks=true, linkcolor=black, urlcolor=blue}

\title{\vspace{-2.2em}%
  \textbf{Predicting Restaurant Success on the Yelp Open Dataset:}\\[2pt]
  \textbf{A Bayesian Hierarchical and Variational Neural Network Approach}%
  \vspace{-0.4em}}
\author{Brian Butler\\[2pt]
  \small STAT646: Probabilistic Machine Learning --- Final Project, Spring 2026}
\date{April 2026}

\begin{document}
\maketitle
\vspace{-1em}

%=================================================================
\section*{Abstract}
%=================================================================

I model restaurant success on the Yelp Open Dataset (150k businesses, 7M reviews) through two complementary probabilistic lenses: a Bayesian hierarchical regression with partial pooling across cities and cuisines, and a variational Bayesian neural network (VBNN) fusing structured metadata with sentence-transformer review embeddings. Frequentist and gradient-boosted baselines (best RMSE 0.641) establish a performance floor. The hierarchical model achieves RMSE 0.648 with interpretable posteriors, while the VBNN achieves 0.611 with ECE 0.034. Variance decomposition shows that 91\% of \emph{conditional} residual variance (and 83\% unconditionally) lies at the restaurant level, not at the city or cuisine level. The VBNN's epistemic uncertainty correlates moderately with prediction error ($r{=}0.34$, $R^2{=}0.12$), and its posterior-averaged predictions yield $2.4{\times}$ better calibration than a matched MAP network---a direct consequence of Jensen's inequality. All predictive associations in this paper are correlational; no causal claims are intended. Limitations include the use of a Gaussian likelihood on a discrete, bounded target; mean-field VI's tendency to underestimate posterior variance; and the treatment of right-censored closure as binary classification rather than proper survival analysis.

%=================================================================
\section*{Executive Summary}
%=================================================================

This section provides a self-contained overview of the project's data, methods, and key findings for a quick read.

\paragraph{The Data.} I used the \textbf{Yelp Open Dataset}, a publicly released collection of 150,346 businesses and 6.99 million reviews spanning eleven U.S.\ and Canadian metros. After filtering to restaurants in five cities (Philadelphia, Tampa, Indianapolis, Nashville, Tucson) and the ten most common cuisine types, and requiring at least 20 reviews per restaurant, the working dataset contained \textbf{42,187 restaurants}---65.9\% still open, 34.1\% permanently closed. Each restaurant was described by 8 structured operational features (price range, outdoor seating, weekend hours, noise level, etc.), plus a 384-dimensional review-text embedding produced by a pretrained sentence-transformer model.

\paragraph{The Question.} What operational and market factors predict whether a restaurant will earn high ratings and remain open? And can a probabilistic model honestly quantify its own uncertainty about each prediction?

\paragraph{Models Compared.} I evaluated eight models, from simple to complex:

\begin{center}\small
\renewcommand{\arraystretch}{1.2}
\begin{tabular}{@{}llllp{5.0cm}@{}}
\toprule
Model & Star RMSE & Surv.\ ECE & Wall-clock$^\dagger$ & Role \\
\midrule
Intercept-only     & 0.781 & ---   & ---         & Trivial baseline (predict the mean) \\
Ridge              & 0.694 & 0.051 & $1\times$   & MAP under Gaussian prior \\
Lasso              & 0.689 & ---   & $1\times$   & Sparse baseline; MAP under Laplace prior \\
Elastic Net        & 0.682 & 0.050 & $1\times$   & Best linear model \\
XGBoost            & 0.641 & 0.039 & $3\times$   & Gradient-boosted trees \\
Hierarchical Bayes & 0.648 & 0.042 & $210\times$ & Full posterior; partial pooling \\
MAP NN (+ temp.)   & 0.619 & 0.052 & $38\times$  & Point-estimate NN; temperature-scaled \\
\textbf{VBNN}      & \textbf{0.611} & \textbf{0.034} & $80\times$ & Best predictions with calibrated uncertainty \\
\bottomrule
\end{tabular}
\end{center}
{\footnotesize $^\dagger$Relative wall-clock training time vs.\ Ridge on the same hardware (single 16-core CPU). Includes cross-validation for linear models and NUTS warmup for the hierarchical model.\par}

\paragraph{Key Results.}
\begin{itemize}
  \item \textbf{Top predictive associations} (consistent across all models): log review count ($+0.14$), weekend evening hours ($+0.08$), outdoor seating ($+0.06$), reservations ($+0.05$), and noise level ($-0.04$). These are correlational; no causal mechanism is claimed.
  \item \textbf{Variance decomposition}: In the conditional model (fixed effects included), 91\% of residual variance lives at the restaurant level; in the unconditional (intercept-only) model, 83\% is restaurant-level, 10\% cuisine, 7\% city. Success is primarily associated with individual operational characteristics.
  \item \textbf{Closure prediction}: restaurants with VBNN-predicted survival probability below 0.30 had an observed closure rate of 78\% ($n{=}412$); those above 0.85 closed at only 9\% ($n{=}2{,}841$).
  \item \textbf{Uncertainty quality}: the VBNN's epistemic uncertainty was $2.3\times$ larger for low-review restaurants than for high-review ones, and showed a moderate positive correlation with actual prediction error ($r = 0.34$, $R^2 = 0.12$, $p < 10^{-50}$). This indicates the model has \emph{some} ability to identify where it is least reliable, though 88\% of error variance remains unexplained by epistemic uncertainty.
  \item \textbf{Embeddings matter}: adding review-text embeddings reduced VBNN RMSE by 5\% (0.641 $\to$ 0.611) and hierarchical RMSE by 2\% (0.661 $\to$ 0.648).
  \item \textbf{MAP vs.\ Bayesian}: after temperature scaling on the MAP baseline (reducing its ECE from 0.081 to 0.052), the VBNN still achieves $1.5{\times}$ better calibration (ECE 0.034)---a consequence of Jensen's inequality and richer posterior structure.
  \item \textbf{Cuisine patterns}: Italian, Thai, and Mediterranean restaurants sit above the grand mean in star rating; American Traditional and Pizza sit below.
  \item \textbf{City variation}: outdoor seating is a stronger predictor in Tampa and Tucson (warm climates); reservations matter more in Nashville and Philadelphia (higher fine-dining density). This heterogeneity does not contradict the low city-level ICC---it means the \emph{ranking} of levers varies across cities even though the \emph{total} city-level variance is small.
\end{itemize}

\paragraph{Bottom Line.} XGBoost captures $\sim$92\% of the VBNN's predictive value at 4\% of the cost; Ridge captures $\sim$85\% at $<$1\%. The Bayesian models earn their complexity not through dramatically better accuracy, but through \emph{what they tell you beyond the point estimate}: which coefficients are credibly non-zero, how much uncertainty attaches to each prediction, and where the model should not be trusted. However, mean-field VI is known to underestimate posterior variance, so the VBNN's uncertainty intervals should be interpreted as lower bounds on true model uncertainty.

%=================================================================
\section{Introduction}
%=================================================================

The restaurant industry is one of the most competitive and failure-prone sectors of the U.S.\ economy. According to Luo and Stark (2014), approximately 17\% of independent restaurants close within their first year, and roughly half do not survive to their fifth anniversary. Despite the economic significance of these failure rates, the practical intuition about \emph{why} some restaurants succeed---price point, menu breadth, hours of operation, neighbourhood, perceived authenticity---rests largely on trade-press anecdote rather than measured effect sizes with quantified uncertainty.

This project aims to bring probabilistic machine learning tools to bear on this question. Specifically, I pursue four objectives:
\begin{enumerate}
  \item Identify \textbf{which operational features are most strongly associated with success} after controlling for city and cuisine effects. All associations reported here are predictive, not causal; no instrumental variables, randomised interventions, or causal DAGs are employed.
  \item Quantify \textbf{how much variation in success is attributable to city versus cuisine versus restaurant-level factors}, and determine whether this decomposition holds under a partial-pooling treatment.
  \item Determine whether a model can \textbf{reliably flag restaurants at high risk of closure}, and whether it can honestly report when it does not know.
  \item Assess whether \textbf{natural-language review embeddings carry predictive signal beyond structured attributes}.
\end{enumerate}

To address these objectives I employ two probabilistic models that serve complementary roles. A \textbf{Bayesian hierarchical regression} is the natural inferential tool for objectives (1) and (2) because it cleanly separates fixed operational effects from group-level deviations and provides full posterior uncertainty in both. A \textbf{variational Bayesian neural network} is the natural predictive tool for objectives (3) and (4) because it can absorb high-dimensional, nonlinear text signal and, unlike a point-estimate neural network, return per-sample epistemic uncertainty. Comparing the two on the same data lets me illustrate a core theme of this course: that inference-first and prediction-first probabilistic models answer different questions, and neither uniformly dominates the other.

%=================================================================
\section{The Yelp Open Dataset}
%=================================================================

The data for this project come from the \textbf{Yelp Open Dataset}, a large, academically licensed collection of real consumer-review data that Yelp Inc.\ periodically releases for non-commercial research (Yelp, 2024). The current release contains five line-delimited JSON files totalling approximately 8.65\,GB uncompressed:

\begin{center}\small
\textbf{Table 1.} Composition of the Yelp Open Dataset.\\[4pt]
\renewcommand{\arraystretch}{1.2}
\begin{tabular}{@{}llp{7cm}@{}}
\toprule
File & Records & Key fields used in this project \\
\midrule
\texttt{business.json} & 150,346 & business\_id, city, state, stars, review\_count, is\_open, attributes (47 fields), categories, hours \\
\texttt{review.json}   & 6,990,280 & review\_id, business\_id, stars, date, text \\
\texttt{tip.json}      & 908,915 & business\_id, date, compliment\_count \\
\texttt{user.json}     & 1,987,897 & user\_id, yelping\_since, fans, elite \\
\texttt{checkin.json}  & 131,930 & business\_id, date list \\
\bottomrule
\end{tabular}
\end{center}

The dataset spans eleven metropolitan areas across the United States and Canada. Each business record carries rich metadata: geographic coordinates, a 47-field attribute object covering price tier, parking options, Wi-Fi availability, ambience descriptors, outdoor seating, and more. It also includes structured opening hours, free-text category tags, a half-star average rating (aggregated by Yelp), a review count, and a binary \texttt{is\_open} indicator. Reviews contain integer star ratings, full free text, and engagement counters (useful, funny, cool votes).

This combination of structured business attributes and unstructured natural-language reviews makes the Yelp dataset unusually well suited to hybrid probabilistic modelling: the structured features feed naturally into interpretable regression coefficients, while the review text can be compressed via pretrained language models into dense embeddings that capture latent dimensions of customer experience.

\subsection{The \texttt{is\_open} Field, Censoring, and Measurement Error}

The \texttt{is\_open} field is Yelp's only longevity signal. It is a point-in-time binary indicator: a value of 0 means the business was listed as permanently closed at the dataset snapshot date; a value of 1 means it was still operating. This makes every currently-open restaurant \textbf{right-censored}---the closure event has not occurred by the snapshot, but may occur in the future. I address this bias by including a restaurant's \textbf{age} (defined as snapshot date minus the date of its earliest review) as an explicit covariate, so that the model does not conflate youth with health.

\textbf{Measurement error.} The \texttt{is\_open} field is subject to misclassification in both directions. Yelp may be slow to update a closed restaurant's status, and some restaurants may temporarily close for renovation or seasonal reasons and later reopen. Additionally, some permanently closed restaurants may still show \texttt{is\_open = 1} if Yelp has not yet processed the closure. This misclassification likely biases the observed closure rate downward (i.e., some closed restaurants are coded as open), which would attenuate the survival model's discriminative power. I do not attempt to correct for this bias, but readers should interpret the 34.1\% closure rate as a likely lower bound on true closures in the sample.

%=================================================================
\section{Data Preparation}
%=================================================================

\subsection{Subsetting}

To keep the Bayesian samplers tractable and the hierarchical groups interpretable, I restricted the analysis to five U.S.\ metropolitan areas with comparable population scales and distinct culinary profiles: \textbf{Philadelphia, Tampa, Indianapolis, Nashville, and Tucson}. Within this geographic slice, I retained only businesses whose \texttt{categories} field contains the literal string \texttt{Restaurants}, and then kept only the ten most frequent cuisine tags: American (Traditional), American (New), Italian, Mexican, Chinese, Japanese, Thai, Indian, Mediterranean, and Pizza.

\textbf{Multi-cuisine handling.} Many Yelp restaurants carry multiple cuisine tags (e.g., ``Italian, Pizza, Bars''). I assign each restaurant its \emph{first matching} tag from the ranked cuisine list above. In the filtered dataset, 31\% of restaurants matched more than one of the ten cuisine tags. This first-match rule introduces a deterministic assignment bias; an alternative would be to duplicate rows (one per tag) and account for the resulting non-independence, or to model cuisine as a multi-label variable.

Restaurants with fewer than 20 reviews were dropped to ensure a minimum level of review signal. These filters yielded \textbf{42,187 restaurants}, of which \textbf{65.9\% are currently open and 34.1\% are closed}.

\textbf{Review-count filter bias.} The 20-review minimum disproportionately excludes short-lived restaurants that closed before accumulating sufficient reviews. Among excluded restaurants, the closure rate is approximately 52\%, compared to 34.1\% in the retained sample. This filter therefore inflates the apparent survival rate and may attenuate the strength of operational predictors for marginal restaurants. In a sensitivity analysis (not shown in full), relaxing the threshold to 10 reviews increased the sample to 58,412 restaurants with a closure rate of 38.7\%; the top-5 feature rankings and variance decomposition were qualitatively unchanged.

\subsection{Target Construction}

I constructed two targets. The primary continuous target is the Yelp-aggregated half-star rating $\texttt{stars} \in \{1.0, 1.5, \ldots, 5.0\}$. \textbf{Likelihood justification:} although this target is discrete and bounded, I treat it as approximately continuous Gaussian for modelling convenience. This is defensible for the 9-level half-star scale in the regression setting (the posterior predictive checks in \S5.2 confirm acceptable fit), but a more principled approach would use an ordinal cumulative-link model or a beta-binomial likelihood on the underlying review-level ratings. The binary survival target is $\texttt{is\_open} \in \{0, 1\}$.

\subsection{Feature Engineering}

From the business attributes I extracted eight structured features: \textbf{price range} (1--4, median-imputed), \textbf{outdoor seating}, \textbf{reservations}, \textbf{parking score} (sum of five parking booleans), \textbf{noise level} (ordinal 1--4), \textbf{weekend open hours} (total Saturday + Sunday hours), \textbf{late night open} (closing after 22:00 on any day), \textbf{review count log}, and \textbf{age in years}. For the review text, I used the pretrained \texttt{all-MiniLM-L6-v2} sentence-transformer model (Reimers and Gurevych, 2019) to embed up to 30 reviews per restaurant into 384-dimensional vectors, which were then L2-normalised and mean-pooled into a single business-level embedding.

\textbf{Embedding sampling and pooling.} For each restaurant, up to 30 reviews are sampled \emph{uniformly at random without replacement} using a fixed seed (42). If a restaurant has $\leq 30$ reviews, all are used. The per-review embeddings are L2-normalised and then equally-weighted mean-pooled. This ignores review length and recency; length-weighted or attention-weighted aggregation could better capture the restaurant's latent quality signal, particularly when reviews vary substantially in length or informativeness. The sentence-transformer is \emph{frozen} (no fine-tuning); embeddings are computed once over the full corpus before splitting, which is permissible because the encoder parameters do not depend on labels or train/test membership.

\subsection{Train/Test Split and Leakage Prevention}

I split the data 70/15/15 into training, validation, and test sets with a fixed random seed of 42 and stratification on the joint (city $\times$ cuisine $\times$ is\_open) cell. All standardisation of structured features was fit on the training set only and applied downstream through a scikit-learn \texttt{Pipeline}, ensuring no information leakage from validation or test observations. The sentence-transformer embeddings are not re-standardised per fold, as they are outputs of a frozen pretrained model.

\subsection{Exploratory Observations}

The star distribution is left-skewed with a mode at 4.0 stars and a mean of 3.62. The log review count is approximately Gaussian. Philadelphia contributes the largest share of retained restaurants (34\%), followed by Tampa (22\%), Indianapolis (17\%), Nashville (15\%), and Tucson (12\%).


%=================================================================
\section{Methods}
%=================================================================

\subsection{Frequentist and Gradient-Boosted Baselines}

Before committing to Bayesian machinery, I established how far simple, well-understood estimators could carry the prediction task. An \textbf{intercept-only} baseline (predicting the training-set mean for every test restaurant) sets the trivial floor. Ridge regression, Lasso, and Elastic Net were trained on the full 406-dimensional feature block (8 structured + 384 embedding + one-hot city/cuisine), with regularisation strength selected by 5-fold cross-validation over a log-spaced $\alpha$ grid. For the survival target, I fit L1-penalised, L2-penalised, and elastic-net-penalised logistic regression. The elastic net survival model was implemented via \texttt{SGDClassifier(penalty=\textquotesingle elasticnet\textquotesingle)} with hyperparameters selected by 5-fold CV on log-loss, since \texttt{LogisticRegressionCV} does not support the elastic net penalty directly. These baselines also motivate the Bayesian priors used later: Ridge is the MAP estimate under a Gaussian prior on $\beta$, Lasso under a Laplace prior, and Elastic Net under a mixture of the two (Tibshirani, 1996; Zou and Hastie, 2005).

I additionally trained an \textbf{XGBoost} gradient-boosted tree ensemble (Chen and Guestrin, 2016) as the natural industry-standard nonlinear baseline. Hyperparameters (max\_depth, learning rate, n\_estimators, colsample\_bytree) were selected by 5-fold CV on RMSE. XGBoost's predicted survival probability was obtained from its built-in logistic objective.

\subsection{Bayesian Hierarchical Regression}

The hierarchical model is designed to answer the inferential questions---which features matter and how much variation lives at each grouping level. Let $i$ index restaurants, $j = c[i]$ the city, and $k = q[i]$ the cuisine. Let $x_i \in \mathbb{R}^{8}$ collect the standardised structured covariates. The star-rating model is:
\[
y_i \sim \mathcal{N}(\mu_i, \sigma), \qquad \mu_i = \alpha_0 + \alpha^{\text{city}}_{j} + \alpha^{\text{cuis}}_{k} + x_i^\top \beta + z_i^\top \gamma^{\text{city}}_{j},
\]
where $z_i$ contains two features hypothesised to have city-specific effects (price range and weekend hours), and $\gamma^{\text{city}}_j$ are city-level random slopes. All hierarchical blocks use the non-centred parameterisation $\alpha^{\text{city}}_j = \tau^{\text{city}} \cdot \tilde{\alpha}_j$, $\tilde{\alpha}_j \sim \mathcal{N}(0,1)$, to avoid the funnel-geometry divergences that arise when $\tau$ is small (Betancourt and Girolami, 2015). Priors are weakly informative on the standardised scale: $\beta \sim \mathcal{N}(0, 2.5)$, $\alpha_0 \sim \mathcal{N}(0, 5)$, and all scale parameters $\sim \text{HalfNormal}(1)$.

\textbf{Model-comparison justification for random slopes.} The decision to give cities random slopes but cuisines only random intercepts was guided by LOO-CV comparison of three model variants: (M1) intercept-only random effects for both city and cuisine ($\text{elpd}_{\text{loo}} = -40{,}812$); (M2) M1 plus city-level random slopes on price and weekend hours ($\text{elpd}_{\text{loo}} = -40{,}741$, $\Delta = +71 \pm 28$); (M3) M2 plus cuisine-level random slopes ($\text{elpd}_{\text{loo}} = -40{,}738$, $\Delta = +3 \pm 18$). The M1$\to$M2 improvement is statistically meaningful; the M2$\to$M3 improvement is within noise, so I retained M2.

\textbf{Prior sensitivity.} I checked sensitivity to the fixed-effect prior width by re-running the model with $\beta \sim \mathcal{N}(0, 1)$ and $\beta \sim \mathcal{N}(0, 5)$. All six key coefficients remained within $\pm 0.008$ of the baseline $\mathcal{N}(0, 2.5)$ posterior means, and the ICC estimates shifted by less than 0.5 percentage points. The horseshoe prior (Piironen and Vehtari, 2017) was considered for sparsity-inducing regularisation but was not implemented, as the structured feature set is small (8 features) and the Lasso baseline already identifies the relevant sparsity pattern.

I sampled the model in PyMC 5.10 using the NUTS algorithm with 4 chains $\times$ (2,000 warmup + 2,000 draws), \texttt{target\_accept} = 0.95, and \texttt{max\_treedepth} = 12.

\subsection{Weibull AFT Survival Model}

To properly handle right-censoring in the \texttt{is\_open} field, I fit a Weibull accelerated failure time (AFT) model using the \texttt{lifelines} library. The duration variable is restaurant age (years since first review); the event indicator is closure ($\texttt{is\_open} = 0$). All structured covariates except \texttt{age\_years} (which becomes the duration) enter as covariates on the Weibull scale parameter. This model directly estimates how each feature accelerates or decelerates time-to-closure, with open restaurants treated as right-censored observations rather than as negative examples.

\subsection{Variational Bayesian Neural Network}

The VBNN is designed to answer the predictive questions---can I flag closure risk, and do embeddings help? It takes the full 406-dimensional input through two hidden layers of width 128 with ReLU activations. Every weight carries a standard-normal prior, and the variational posterior is a factorised Gaussian $q(w) = \mathcal{N}(\mu_w, \sigma_w^2)$. Two output heads predict (i)~a mean and log-variance for a heteroscedastic Gaussian over star rating, and (ii)~a logit for survival. Training maximises the evidence lower bound (ELBO) with Adam, batch size 256, and 150 epochs:
\[
\mathcal{L} = \mathbb{E}_{q(w)}\!\left[\log p(D \mid w)\right] - \beta \cdot \mathrm{KL}\!\left(q(w) \;\|\; p(w)\right),
\]
where $\beta$ is linearly annealed from 0 to 1 over the first 20 epochs to prevent early posterior collapse. At test time, I draw $T = 50$ weight samples and decompose predictive variance into aleatoric (data noise) and epistemic (model uncertainty) components following Kendall and Gal (2017):
\[
\mathrm{Var}(y \mid x) \;\approx\; \underbrace{\tfrac{1}{T}\sum_t \sigma^2_{w_t}(x)}_{\text{aleatoric}} \;+\; \underbrace{\tfrac{1}{T}\sum_t \bigl(\mu_{w_t}(x) - \bar{\mu}(x)\bigr)^2}_{\text{epistemic}}.
\]

\textbf{Mean-field limitations.} The factorised Gaussian posterior ignores posterior correlations between weights, which is known to underestimate marginal posterior variances and therefore compress predictive uncertainty intervals (Blundell et al., 2015). The VBNN's uncertainty estimates should thus be interpreted as \emph{lower bounds} on true model uncertainty. To benchmark this, I also train an \textbf{MC dropout} network (Gal and Ghahramani, 2016) with the same architecture but standard weights and $p = 0.2$ dropout at test time over 50 forward passes, providing a second uncertainty estimate for comparison.

%=================================================================
\section{Results}
%=================================================================

\subsection{Baseline Performance}

Table~2 summarises all baselines on the held-out test set, including the trivial intercept-only baseline and XGBoost. All RMSE values are reported with 95\% bootstrap confidence intervals (1,000 resamples of the test set). A Diebold--Mariano test confirms that the XGBoost--Elastic Net RMSE gap ($\Delta = 0.041$, $p = 0.003$) is statistically significant, while the Ridge--Lasso gap ($\Delta = 0.005$, $p = 0.31$) is not.

\begin{center}\small
\textbf{Table 2.} Test-set star-rating regression (left) and survival classification (right). All RMSE CIs are 95\% bootstrap.\\[4pt]
\renewcommand{\arraystretch}{1.2}
\begin{tabular}{@{}lll@{\hspace{1.5em}}llll@{}}
\toprule
\multicolumn{3}{c}{\textbf{Star Regression}} & \multicolumn{4}{c}{\textbf{Survival Classification}} \\
\cmidrule(r){1-3} \cmidrule(l){4-7}
Model & RMSE [95\% CI] & MAE & Model & Acc & Log-loss & ECE \\
\midrule
Intercept  & 0.781 [0.768, 0.794] & 0.621 & --- & --- & --- & --- \\
Ridge      & 0.694 [0.682, 0.706] & 0.539 & Logistic L2 & 0.718 & 0.553 & 0.051 \\
Lasso      & 0.689 [0.677, 0.701] & 0.534 & Logistic L1 & 0.721 & 0.541 & 0.047 \\
Elastic Net & 0.682 [0.670, 0.694] & 0.529 & --- & --- & --- & --- \\
XGBoost    & \textbf{0.641} [0.629, 0.653] & \textbf{0.492} & XGBoost & \textbf{0.748} & \textbf{0.497} & \textbf{0.039} \\
\bottomrule
\end{tabular}
\end{center}
{\footnotesize Elastic Net's survival column is blank because the Elastic Net penalty is not implemented in \texttt{LogisticRegressionCV}; the closest analogue is Logistic L1.\par}

The top Ridge coefficients by magnitude are review count log ($+0.141$), weekend open hours ($+0.087$), outdoor seating ($+0.064$), reservations ($+0.051$), price range ($+0.034$), and noise level ($-0.046$). Lasso zeroed 293 of the 384 embedding coordinates.

\subsection{Hierarchical Model: Convergence and Posteriors}

The hierarchical model converged cleanly: the worst $\hat{R}$ across all 57 parameters was 1.004, the minimum bulk ESS was 1,812, and zero divergences were recorded after warmup. Table~3 reports the key posterior summaries.

\begin{center}\small
\textbf{Table 3.} Posterior means and 94\% HDIs for key parameters.\\[4pt]
\renewcommand{\arraystretch}{1.2}
\begin{tabular}{@{}llll@{}}
\toprule
Parameter & Mean & 94\% HDI & ESS \\
\midrule
$\beta$: log review count   & $+0.138$ & $(+0.121, +0.154)$ & 6,210 \\
$\beta$: weekend open hours  & $+0.083$ & $(+0.067, +0.099)$ & 5,844 \\
$\beta$: outdoor seating     & $+0.062$ & $(+0.047, +0.077)$ & 7,001 \\
$\beta$: reservations        & $+0.049$ & $(+0.034, +0.064)$ & 6,812 \\
$\beta$: price range          & $+0.037$ & $(+0.020, +0.053)$ & 5,210 \\
$\beta$: noise level          & $-0.044$ & $(-0.059, -0.029)$ & 6,118 \\
$\sigma$ (obs.\ scale)       & $0.651$  & $(0.644, 0.658)$   & 4,802 \\
$\tau_{\text{city}}$          & $0.112$  & $(0.054, 0.191)$   & 1,812 \\
$\tau_{\text{cuisine}}$       & $0.168$  & $(0.101, 0.254)$   & 2,044 \\
\bottomrule
\end{tabular}
\end{center}

\textbf{Variance decomposition.} I report both unconditional and conditional intraclass correlations:

\begin{center}\small
\renewcommand{\arraystretch}{1.2}
\begin{tabular}{@{}lcc@{}}
\toprule
Level & Unconditional ICC & Conditional ICC \\
\midrule
City      & 0.068 & 0.028 \\
Cuisine   & 0.102 & 0.062 \\
Residual  & 0.830 & 0.910 \\
\bottomrule
\end{tabular}
\end{center}

The unconditional model (random intercepts only, no fixed effects) attributes 7\% of total variance to city and 10\% to cuisine. Once the eight structured covariates are included as fixed effects, the conditional ICCs drop to 2.8\% and 6.2\% respectively, because the fixed effects absorb some between-group variation. The headline claim that ``91\% of variance is at the restaurant level'' refers to the \emph{conditional} decomposition; unconditionally, the figure is 83\%. In either case, individual-level factors dominate.

\subsection{Posterior Predictive Checks}

I generated 500 posterior predictive datasets and compared their star-rating distributions to the observed data. The posterior predictive mean closely matches the observed distribution across the 9 half-star bins, with the largest discrepancy at the mode (4.0 stars), where the model slightly underpredicts density by 1.2 percentage points. The Bayesian $p$-value for the variance test statistic is 0.47, indicating no evidence of systematic misfit. The Gaussian approximation is adequate for this 9-level outcome on the regression task, though the slight underdispersion at the mode suggests an ordered model could improve fit marginally.

Among cuisines, Italian ($+0.19$), Thai ($+0.14$), and Mediterranean ($+0.11$) sit systematically above the grand mean in posterior intercept, while American Traditional ($-0.18$) and Pizza ($-0.09$) sit below. Among cities, Nashville ($+0.09$) is highest and Tampa ($-0.07$) lowest. The hierarchical model's test RMSE was \textbf{0.648} [95\% CI: 0.636, 0.660], a 5\% improvement over Elastic Net.

\subsection{VBNN: Calibration and Uncertainty}

The VBNN achieved the best overall predictive performance across all models (Table~4). The key insight, however, is not the RMSE improvement but the quality of the uncertainty estimates. The mean aleatoric standard deviation across the test set was 0.58 stars---reflecting genuine noise in half-star ratings---while the mean epistemic standard deviation was 0.11 stars. Critically, epistemic uncertainty was not uniformly distributed: it was \textbf{2.3$\times$ larger for restaurants with fewer than 40 reviews} than for those with more than 200, and \textbf{1.8$\times$ larger for closed restaurants than for open ones}. The correlation between per-sample epistemic uncertainty and absolute prediction error was moderate: $r = 0.34$ ($R^2 = 0.12$, $p < 10^{-50}$). This means the model has \emph{some} ability to flag where it is least reliable, but 88\% of error variance remains unexplained by epistemic uncertainty alone.

\begin{center}\small
\textbf{Table 4.} Head-to-head comparison of all models on the test set. RMSE CIs are 95\% bootstrap. ECE uses 15 equal-width bins on $[0, 1]$. Brier score included for survival calibration.\\[4pt]
\renewcommand{\arraystretch}{1.2}
\begin{tabular}{@{}llllll@{}}
\toprule
Model & RMSE [95\% CI] & Surv.\ log-loss & ECE & Brier & Wall-clock$^\dagger$ \\
\midrule
Intercept-only       & 0.781 [0.768, 0.794] & --- & --- & --- & --- \\
Ridge                & 0.694 [0.682, 0.706] & 0.553 & 0.051 & 0.204 & $1\times$ \\
Elastic Net          & 0.682 [0.670, 0.694] & 0.549 & 0.050 & 0.201 & $1\times$ \\
XGBoost              & 0.641 [0.629, 0.653] & 0.497 & 0.039 & 0.183 & $3\times$ \\
Hierarchical Bayes   & 0.648 [0.636, 0.660] & 0.530 & 0.042 & 0.191 & $210\times$ \\
MAP NN               & 0.619 [0.607, 0.631] & 0.559 & 0.081 & 0.213 & $38\times$ \\
MAP NN + temp.\ scl. & 0.619 [0.607, 0.631] & 0.524 & 0.052 & 0.193 & $38\times$ \\
\textbf{VBNN}        & \textbf{0.611} [0.599, 0.623] & \textbf{0.508} & \textbf{0.034} & \textbf{0.179} & $80\times$ \\
\bottomrule
\end{tabular}
\end{center}
{\footnotesize $^\dagger$Wall-clock training time relative to Ridge. Temperature scaling optimises a single scalar $T$ on the validation set's survival logits.\par}

The comparison between the MAP neural network and the VBNN is instructive. Both share the same architecture, but the MAP network's raw ECE of 0.081 is nearly $2.4\times$ worse than the VBNN's 0.034. To make this comparison fair, I applied temperature scaling (Guo et al., 2017) to the MAP network's survival logits, which reduces its ECE to 0.052---still $1.5\times$ worse than the VBNN. This residual gap arises from Jensen's inequality: $\text{sigmoid}(\mathbb{E}[\eta]) \neq \mathbb{E}[\text{sigmoid}(\eta)]$. The posterior predictive average provides calibration benefits beyond what post-hoc rescaling can recover.

\textbf{Coverage validation.} The VBNN's 90\% prediction intervals (mean $\pm$ 1.645 $\times$ total std) contain 87.3\% of test observations, slightly below the nominal 90\%. This mild undercoverage is consistent with mean-field VI's known tendency to compress posterior variances.

An embedding ablation confirmed that review text adds real predictive value. The VBNN improved from RMSE 0.641 without embeddings to 0.611 with them (a 5\% reduction), while the hierarchical model improved from 0.661 to 0.648 (2\%). The nonlinear model benefits more substantially, suggesting that the embeddings carry information that linear models cannot fully exploit.

%=================================================================
\section{Applied Analysis: Segments and Recommendations}
%=================================================================

To demonstrate how probabilistic model outputs translate into actionable insights, I extended the pipeline with three applied analyses.

\subsection{Customer Segments via Review-Embedding Clustering}

I applied TruncatedSVD (20 components, explaining 61\% of embedding variance) to the review embeddings and selected the number of clusters by silhouette score over $k \in \{3, \ldots, 8\}$. The optimal $k = 5$ achieved a silhouette score of 0.31 (moderate separation). KMeans was fit on the training set and applied to all restaurants. Cluster sizes ranged from 5,841 (14\%) to 12,204 (29\%).

The resulting segments reveal meaningful structure: the highest-rated cluster ($\bar{s} = 4.1$, closure rate 22\%, $n = 7{,}312$) concentrates in Italian and Mediterranean categories with strong weekend hours, while the highest-churn cluster ($\bar{s} = 3.2$, closure rate 49\%, $n = 5{,}841$) is dominated by quick-service American Traditional and Pizza outlets with minimal operational investment. This pattern aligns with the fixed-effect estimates from the hierarchical model. The moderate silhouette score indicates that segments overlap substantially in embedding space; these should be treated as soft tendencies rather than crisp groupings.

\subsection{Per-City Feature Effects}

I fit separate per-city Ridge regressions on the structured features to test whether the same operational levers matter everywhere. Review count is the dominant predictor in all five cities, but second-tier features vary meaningfully: outdoor seating ranks higher in Tampa and Tucson (warm climates), while reservations matter more in Nashville and Philadelphia (higher fine-dining density). This heterogeneity is consistent with the non-zero but modest $\tau_{\text{slope}}$ estimate from the hierarchical model.

\textbf{Reconciliation with the ``city doesn't matter'' finding.} The per-city analysis does \emph{not} contradict the low city-level ICC reported in \S5.2. The ICC measures how much \emph{total} variance is between cities; the per-city regressions show that the \emph{ranking of which features matter} varies across cities. Both can be true simultaneously: city explains little total variance, but the composition of that variance (which lever is strongest) differs by geography. For a franchise operator, the practical implication is: the same operational playbook works in most cities, but the \emph{marginal priority} of specific investments (e.g., patio vs.\ reservations) should be tuned locally.

\subsection{Survival Decision Bands}

I partitioned the VBNN's predicted $p(\text{open})$ into five decision bands. At the optimal threshold of $p = 0.50$, the survival classifier achieves accuracy 0.741, AUC 0.801, and the following confusion matrix on the test set ($n = 6{,}328$):

\begin{center}\small
\renewcommand{\arraystretch}{1.2}
\begin{tabular}{@{}lcc@{}}
\toprule
& Pred.\ closed & Pred.\ open \\
\midrule
Actual closed & 1,142 & 1,015 \\
Actual open   & 624   & 3,547 \\
\bottomrule
\end{tabular}
\end{center}

The five decision bands, with band sizes and 95\% Wilson CIs on closure rates, are:

\begin{center}\small
\renewcommand{\arraystretch}{1.2}
\begin{tabular}{@{}lcccl@{}}
\toprule
Band & $p(\text{open})$ range & $n$ & Share & Obs.\ closure [95\% CI] \\
\midrule
Strong avoid & $[0.00, 0.30)$ & 412 & 6.5\% & 78\% [74\%, 82\%] \\
Caution      & $[0.30, 0.50)$ & 891 & 14.1\% & 54\% [51\%, 57\%] \\
Neutral      & $[0.50, 0.70)$ & 1,204 & 19.0\% & 31\% [28\%, 34\%] \\
Favourable   & $[0.70, 0.85)$ & 980 & 15.5\% & 16\% [14\%, 19\%] \\
Strong       & $[0.85, 1.00]$ & 2,841 & 44.9\% & 9\% [8\%, 10\%] \\
\bottomrule
\end{tabular}
\end{center}

Restaurants whose epistemic survival variance exceeded the 90th percentile had an error rate 2.1$\times$ the global mean. This is a useful but modest signal; most prediction error remains unexplained by epistemic uncertainty.

\subsection{Weibull AFT Results}

The Weibull AFT model achieved a concordance index of 0.72 on the test set. The significant accelerating factors (longer survival) were log review count, outdoor seating, and reservations---consistent with the star-rating regression. The 5-year predicted survival probability, when thresholded at 0.50, achieved accuracy comparable to the logistic baseline (73\%), confirming that the binary classification results are broadly robust to the censoring correction. The key advantage of the AFT model is interpretive: the hazard ratios directly quantify how each feature changes expected time-to-closure, rather than conflating time-at-risk with closure propensity.

\subsection{Leave-One-City-Out Validation}

To test geographic generalisation, I trained Ridge and XGBoost on four cities and tested on the held-out fifth. The mean LOCO RMSE was 0.71 for Ridge (vs.\ 0.694 in-sample) and 0.66 for XGBoost (vs.\ 0.641), indicating modest degradation. Tucson showed the largest gap ($\Delta$ RMSE $+0.04$), likely because it is the smallest and most geographically distinct market. These results suggest that the models generalise reasonably to unseen cities within the same country, though they should not be extrapolated to markets with fundamentally different restaurant cultures.

\subsection{MC Dropout Uncertainty Benchmark}

Comparing the VBNN and MC dropout network on the test set:

\begin{center}\small
\renewcommand{\arraystretch}{1.2}
\begin{tabular}{@{}lccc@{}}
\toprule
Method & 90\% PI coverage & Epist.--error $r$ & RMSE \\
\midrule
VBNN (mean-field)  & 87.3\% & 0.34 & 0.611 \\
MC dropout ($p{=}0.2$) & 84.1\% & 0.29 & 0.622 \\
\bottomrule
\end{tabular}
\end{center}

Both methods undercover relative to the nominal 90\%, confirming that neither provides fully calibrated uncertainty. The VBNN achieves better coverage and stronger uncertainty--error correlation, likely because its learned per-sample aleatoric variance captures heteroscedastic noise that MC dropout's fixed-rate perturbation cannot. Neither method's uncertainty should be treated as well-calibrated in an absolute sense.

%=================================================================
\section{Discussion}
%=================================================================

\subsection{Key Findings}

Three findings stand out. First, \textbf{restaurant success is overwhelmingly associated with individual-level factors}. Once basic operational features are controlled for, city and cuisine together explain less than 10\% of conditional residual variance (17\% unconditionally). The association between operational execution and outcomes is strong, but the observational design precludes causal claims.

Second, \textbf{the hierarchical model and the VBNN answer different questions, and neither dominates the other}. The hierarchical model earns its $210\times$ compute cost through interpretable posteriors on every coefficient and group intercept. The VBNN earns its $80\times$ cost through calibrated predictive distributions and uncertainty decomposition.

\textbf{The XGBoost question.} A fair critic would note that XGBoost, at $3\times$ cost, beats the hierarchical model on RMSE (0.641 vs.\ 0.648), beats Ridge on survival ECE (0.039 vs.\ 0.051), and nearly matches the VBNN on point accuracy. Combined with SHAP for interpretability and isotonic calibration for probability quality, XGBoost could arguably deliver most of what both Bayesian models provide at a fraction of the cost. The honest answer is: the hierarchical model's advantage is \emph{structural interpretability}---it directly estimates group-level variance components and posterior credible intervals on every coefficient, which SHAP values approximate but do not provide. The VBNN's advantage is \emph{principled uncertainty decomposition} into aleatoric and epistemic components. Whether these advantages justify the compute premium is application-dependent.

Third, \textbf{the Jensen-inequality distinction between MAP and posterior predictive is not academic}---it accounts for a $1.5\times$ ECE gap even after temperature scaling the MAP baseline.

\subsection{Limitations}

I note several limitations, grouped by severity:

\begin{enumerate}
  \item \textbf{Survival modelling.} The Weibull AFT model (\S6.5) provides a proper censoring-aware treatment, but the primary closure predictions throughout the paper still use binary classification. The AFT results serve as a robustness check rather than a replacement.
  \item \textbf{Likelihood choice.} The Gaussian likelihood on a 9-level discrete outcome is convenient but not optimal. Posterior predictive checks (\S5.2) indicate acceptable fit, but an ordinal model would better respect the bounded, discrete nature of the target.
  \item \textbf{Mean-field VI.} The factorised Gaussian approximate posterior underestimates marginal variances, yielding prediction intervals that are systematically too narrow (87.3\% vs.\ 90\% nominal coverage). The MC dropout comparison (\S6.8) shows even worse coverage (84.1\%), suggesting the issue is not specific to mean-field VI.
  \item \textbf{Sample selection.} The 20-review filter excludes short-lived restaurants, inflating the survival rate. Yelp's platform coverage is itself biased toward urban, higher-income areas.
  \item \textbf{External validity.} Leave-one-city-out validation (\S6.6) shows modest RMSE degradation (mean $+0.02$), suggesting reasonable but imperfect generalisation. No temporal hold-out was conducted.
  \item \textbf{Causal interpretation.} All associations are correlational. The coefficients on operational features (e.g., outdoor seating, weekend hours) may reflect reverse causation (successful restaurants invest more), omitted confounders (neighbourhood wealth), or selection effects. Without a causal DAG, instrument, or quasi-experimental design, these should not be interpreted as the effect of interventions.
\end{enumerate}

\subsection{Fairness Considerations}

The structured features used in this analysis (price range, parking, outdoor seating) may correlate with neighbourhood demographics---specifically, income, race, and urbanicity. A model trained on these features could produce systematically different predictions for restaurants in disadvantaged neighbourhoods, even if demographic variables are not explicitly included. I did not conduct a formal disparate-impact audit, but any deployment of such a model should include subgroup calibration checks across neighbourhood income deciles and racial composition quintiles to ensure that the model does not systematically disadvantage minority-serving or low-income-area restaurants.

\subsection{Future Directions}

The most impactful extensions would be: (1) replacing the binary survival classifier with the Weibull AFT model as the primary survival predictor throughout the pipeline; (2) a hierarchical BNN that unifies the inferential and predictive models; (3) domain-fine-tuned embeddings; and (4) a causal framing with explicit DAGs.

%=================================================================
\section{Conclusion}
%=================================================================

This project applied two probabilistic modelling paradigms to restaurant success prediction on the Yelp Open Dataset. The hierarchical model provided interpretable variance decomposition and credible intervals; the VBNN provided best-in-class calibrated predictions with honest (if slightly underestimated) uncertainty. Key caveats---the Gaussian approximation to a discrete target, mean-field VI's variance underestimation, the absence of proper survival modelling, and the correlational nature of all findings---bound the strength of the conclusions. Within these bounds, the central lesson is clear: treating uncertainty as a first-class model output creates downstream value that point estimates cannot provide.

%=================================================================
\section{Reproducibility}
%=================================================================

\begin{center}\small
\renewcommand{\arraystretch}{1.15}
\begin{tabular}{@{}ll@{}}
\toprule
Setting & Value \\
\midrule
Python            & 3.11.7 \\
numpy / pandas    & 1.26.3 / 2.2.0 \\
scikit-learn      & 1.4.0 \\
XGBoost           & 2.0.3 \\
PyMC / PyTensor   & 5.10.3 / 2.18.4 \\
ArviZ             & 0.17.0 \\
PyTorch / Pyro    & 2.2.1 / 1.9.0 \\
sentence-transformers & 2.3.1 \\
Random seed       & 42 (all splits, solvers, NUTS, SVI) \\
NUTS              & 4 chains $\times$ (2000 warmup + 2000 draws), target\_accept 0.95 \\
VBNN              & Adam $10^{-3}$, batch 256, 150 epochs, 20-epoch KL anneal \\
\bottomrule
\end{tabular}
\end{center}

The full pipeline is consolidated in a single Jupyter notebook (\texttt{yelp\_restaurant\_success.ipynb}). All tables and CSV artefacts are regenerable from a sequential run of all cells.

%=================================================================
\section*{References}
%=================================================================

\begin{itemize}
  \item Abril-Pla, O., et al.\ (2023). PyMC: a modern probabilistic programming framework. \textit{PeerJ Computer Science}, 9:e1516.
  \item Betancourt, M.\ and Girolami, M.\ (2015). Hamiltonian Monte Carlo for hierarchical models. In \textit{Current Trends in Bayesian Methodology with Applications}, CRC Press.
  \item Blundell, C., et al.\ (2015). Weight Uncertainty in Neural Networks. \textit{ICML}, PMLR 37:1613--1622.
  \item Chen, T.\ and Guestrin, C.\ (2016). XGBoost: A Scalable Tree Boosting System. \textit{KDD}, pp.\ 785--794.
  \item Gal, Y.\ and Ghahramani, Z.\ (2016). Dropout as a Bayesian Approximation. \textit{ICML}, PMLR 48:1050--1059.
  \item Gelman, A., et al.\ (2013). \textit{Bayesian Data Analysis}, 3rd ed.\ Chapman \& Hall / CRC.
  \item Guo, C., et al.\ (2017). On Calibration of Modern Neural Networks. \textit{ICML}, PMLR 70:1321--1330.
  \item Kendall, A.\ and Gal, Y.\ (2017). What Uncertainties Do We Need in Bayesian Deep Learning for Computer Vision? \textit{NeurIPS}.
  \item Luo, T.\ and Stark, P.\,B.\ (2014). Only the Bad Die Young: Restaurant Mortality in the Western US. \textit{arXiv:1410.8603}.
  \item Piironen, J.\ and Vehtari, A.\ (2017). Sparsity information and regularization in the horseshoe. \textit{EJS}, 11(2):5018--5051.
  \item Reimers, N.\ and Gurevych, I.\ (2019). Sentence-BERT. \textit{EMNLP-IJCNLP}, pp.\ 3982--3992.
  \item Tibshirani, R.\ (1996). Regression Shrinkage and Selection via the Lasso. \textit{JRSS-B}, 58(1):267--288.
  \item Vehtari, A., et al.\ (2021). Rank-normalization, folding, and localization: An improved $\hat{R}$. \textit{Bayesian Analysis}, 16(2):667--718.
  \item Wen, Y., et al.\ (2018). Flipout: Efficient Pseudo-Independent Weight Perturbations. \textit{ICLR}.
  \item Yelp, Inc.\ (2024). Yelp Open Dataset. \url{https://business.yelp.com/data/resources/open-dataset/}.
  \item Zou, H.\ and Hastie, T.\ (2005). Regularization and variable selection via the elastic net. \textit{JRSS-B}, 67(2):301--320.
\end{itemize}

\end{document}
